\chapter{Introdução}

Os sistemas elétricos de potência são estruturas fundamentais para o fornecimento de energia elétrica à sociedade moderna, abrangendo as etapas de geração, transmissão, distribuição e consumo. Para que esses sistemas operem de forma segura e confiável, é necessário empregar dispositivos capazes de identificar condições anormais de operação e atuar de maneira adequada para isolar trechos defeituosos. Nesse contexto, os sistemas de proteção exercem papel essencial, pois são responsáveis por detectar faltas, sobrecorrentes, variações de tensão e outras perturbações que possam comprometer a integridade dos equipamentos e a continuidade do fornecimento de energia \cite{anderson1999power,blackburn2006protective}.

Entre os dispositivos utilizados na proteção de sistemas elétricos de potência, destacam-se os relés de proteção. Esses equipamentos recebem sinais de tensão e corrente do sistema, processam essas grandezas e tomam decisões de atuação com base em critérios previamente definidos. Com o avanço da tecnologia digital, os relés passaram a empregar algoritmos de processamento de sinais, estimação fasorial e lógicas de proteção configuráveis, permitindo maior flexibilidade, precisão e integração com sistemas de supervisão e análise de eventos \cite{phadke2009computer,horowitz2014power}.

Apesar da importância desses conceitos para a formação do engenheiro eletricista, o ensino de proteção de sistemas elétricos de potência ainda é, em muitos casos, predominantemente teórico. A formação de futuros engenheiros de proteção demanda contato com fundamentos, ajustes, testes e análise de desempenho de relés, pois a área envolve não apenas conhecimento conceitual, mas também interpretação de eventos e tomada de decisão baseada em sinais reais ou simulados \cite{brahma2009education}. Entretanto, a utilização de equipamentos reais de proteção, malas de teste e simuladores especializados costuma ser limitada em ambientes acadêmicos devido ao custo elevado, à complexidade de operação e à necessidade de infraestrutura laboratorial específica. Como consequência, temas como atuação de relés, análise de oscilografias, ajustes de proteção, estimação de fasores e comportamento dos sinais durante faltas podem permanecer distantes da experiência prática dos estudantes.

Diante dessa limitação, torna-se relevante o desenvolvimento de ferramentas didáticas acessíveis que permitam aproximar os conteúdos teóricos de uma experiência prática em bancada. Na educação em engenharia, atividades laboratoriais e estratégias de aprendizagem ativa são reconhecidas por favorecer a compreensão de fenômenos físicos, o desenvolvimento de habilidades experimentais e a conexão entre teoria e prática \cite{prince2004active,feisel2005laboratory}. No ensino de proteção de sistemas elétricos, plataformas laboratoriais têm sido propostas justamente para permitir que estudantes configurem relés, observem condições de falta e avaliem mecanismos de atuação de forma experimental \cite{ferris2013protectionlab,radhi2020plc}. Nesse contexto, o uso de microcontroladores de baixo custo, como o ESP32, associado a conversores digital-analógicos e a arquivos padronizados de oscilografia, possibilita a construção de uma plataforma experimental capaz de reproduzir sinais de tensão e corrente e avaliar algoritmos de proteção em condições controladas.

Neste trabalho, propõe-se uma plataforma didática embarcada para reprodução de arquivos COMTRADE e emulação de funções de relés de proteção aplicadas a sistemas elétricos de potência. O padrão COMTRADE é amplamente utilizado para armazenamento e intercâmbio de registros oscilográficos, permitindo representar sinais provenientes de simulações ou medições reais \cite{ieee_comtrade,iec_60255_24_2013}. No projeto desenvolvido, as oscilografias são obtidas a partir de simulações de faltas trifásicas realizadas no Simulink, convertidas para arquivos COMTRADE e posteriormente reproduzidas por uma ESP32 geradora com auxílio de um conversor digital-analógico MCP4728 \cite{espressif_esp32,microchip_mcp4728}.

A plataforma é constituída basicamente por duas unidades ESP32. A primeira unidade é responsável pela geração e reprodução dos sinais analógicos de tensão e corrente a partir das amostras provenientes dos arquivos COMTRADE. A segunda unidade realiza a recepção dos sinais reproduzidos, a aquisição das amostras e a execução dos algoritmos de proteção. Nas versões mais recentes do projeto, o sistema passou a ter maior grau de embarcamento, com a comunicação, a reprodução, a recepção e os algoritmos de proteção concentrados nas ESP32, enquanto os scripts em computador atuam principalmente na preparação dos dados, configuração, supervisão e interface de execução.

Do ponto de vista acadêmico, a principal contribuição deste trabalho está na construção de uma plataforma acessível, reconfigurável e orientada ao ensino, capaz de permitir que estudantes observem a relação entre uma falta simulada, a oscilografia correspondente, a reprodução física dos sinais e a resposta dos algoritmos de proteção. Dessa forma, busca-se complementar o ensino teórico de proteção de sistemas elétricos de potência com uma ferramenta prática, de baixo custo e compatível com a realidade de laboratórios didáticos, reforçando a importância de experiências laboratoriais no processo de formação em engenharia \cite{feisel2005laboratory,ferris2013protectionlab}.

\section{Objetivos}

\subsection{Objetivo geral}

Desenvolver uma plataforma didática embarcada, baseada em microcontroladores ESP32, para reprodução de sinais provenientes de arquivos COMTRADE e emulação de funções de relés de proteção aplicadas a sistemas elétricos de potência.

\subsection{Objetivos específicos}

Para alcançar o objetivo geral, foram definidos os seguintes objetivos específicos:

\begin{itemize}
  \item realizar simulações de faltas trifásicas em um sistema elétrico de potência no ambiente Simulink/Matlab;
  \item exportar as amostras de tensão e corrente obtidas nas simulações para o ambiente de processamento do Matlab;
  \item converter as amostras simuladas para o padrão COMTRADE, gerando os arquivos \texttt{.cfg} e \texttt{.dat};
  \item validar os arquivos COMTRADE quanto ao tipo de arquivo, número de canais, taxa de amostragem e compatibilidade com a plataforma desenvolvida;
  \item desenvolver o sistema embarcado de geração dos sinais analógicos de tensão e corrente utilizando ESP32 e conversor digital-analógico MCP4728;
  \item implementar a comunicação entre o computador e a ESP32 geradora para envio, organização e reprodução das amostras;
  \item desenvolver o sistema embarcado de recepção dos sinais analógicos reproduzidos pela unidade geradora;
  \item processar as amostras recebidas para extração de grandezas utilizadas pelos algoritmos de proteção;
  \item implementar funções de proteção, como sobrecorrente, direcional e distância, conforme os recursos disponíveis na plataforma;
  \item permitir a configuração dos algoritmos de proteção por meio de comandos e parâmetros de ajuste, aproximando a operação do sistema da lógica de configuração de relés reais;
  \item avaliar a resposta da plataforma diante das oscilografias de falta reproduzidas;
  \item analisar a aplicabilidade do sistema como ferramenta didática para o ensino de proteção de sistemas elétricos de potência.
\end{itemize}

\section{Organização do trabalho}

Este trabalho está organizado em seis capítulos. O Capítulo 1 apresenta a introdução, contextualizando o tema, a motivação do projeto, a proposta da plataforma desenvolvida e os objetivos do trabalho.

O Capítulo 2 apresenta a fundamentação teórica necessária ao desenvolvimento do projeto, abordando conceitos relacionados a sistemas elétricos de potência, curtos-circuitos, relés de proteção, arquivos COMTRADE e processamento digital de sinais aplicado à proteção.

O Capítulo 3 descreve a metodologia adotada, incluindo a geração das oscilografias no Simulink, a conversão das amostras para o padrão COMTRADE, a validação dos arquivos e a arquitetura geral da plataforma composta pelas unidades de geração e recepção.

O Capítulo 4 apresenta a implementação do sistema, detalhando os principais componentes de hardware e software, os scripts desenvolvidos, os firmwares das ESP32, o protocolo de comunicação, a bufferização das amostras, a reprodução dos sinais e a estrutura dos algoritmos de proteção.

O Capítulo 5 apresenta os resultados e discussões, contemplando a análise dos sinais reproduzidos e recebidos, a resposta das funções de proteção implementadas, as limitações observadas e a aplicabilidade didática da plataforma.

Por fim, o Capítulo 6 apresenta as considerações finais do trabalho, destacando as principais contribuições, as limitações do sistema desenvolvido e sugestões para trabalhos futuros.
